\chapter[Piecewise-Smooth Functionals]{Piecewise-Smooth Security Functionals}\label{ch:sec:piecewise-smooth}

This chapter extends the geometric theory to non-differentiable functionals, including min-entropy, Bayes vulnerability, and guessing entropy. These are the security measures most relevant to side-channel analysis (see Chapter~\ref{ch:sec:applications-of-the-wis-framework} for detailed applications).


\section{Introduction}\label{ch:sec:introduction-piecewise_smooth}
Chapters~\ref{ch:sec:universal-geometric-bounds}--\ref{ch:sec:the-average-hessian-operator-and-universal-bregman-theo} developed a geometric theory for smooth information
functionals. Many security measures used in information theory and
cryptography are not globally differentiable because they involve
maxima, argmax operators, or probability rankings
\cite{smith2009foundations, malacaria2007assessing,
clark2007static, chatzikokolakis2005probability}.
The key observation is that these non-smoothness phenomena are highly
structured. They occur only when the ordering of posterior probabilities
changes. Between such events, the functionals are ordinary smooth
functions.
This chapter formalizes that observation.
\paragraph{Rank-Based Security Functionals.}
Representative examples include
\begin{itemize}
\item Min-Entropy,
\item Bayes Vulnerability,
\item Guessing Entropy,
\item Success Probability,
\item Top-\(k\) Success Probability,
\item rank-based leakage measures.
\end{itemize}
Each depends on the ordering of posterior probabilities.



\section{Ranking Transition Times}\label{ch:sec:ranking-transition-times}
Let
\[ P(t)=F^{-1}\!\left((1-t)F(P)+tF(Q)\right),
\qquad t\in[0,1], \]
be the WIS geodesic.
Along this geodesic,
$P_i(t)\propto P_i^{1-t}Q_i^{t}.$
A ranking change between coordinates \(i\) and \(j\) occurs only when
$P_i(t)=P_j(t).$
\begin{proposition}[Crossing Time]
\label{ch:prop:17-1}
If
$\log(P_i/P_j)\neq\log(Q_i/Q_j),$
there is at most one crossing time, given by
\[ \boxed{
t_{ij}=
\frac{\log(P_i/P_j)}
{\log(P_i/P_j)-\log(Q_i/Q_j)}.
} \]
Only values satisfying \(t_{ij}\in(0,1)\) correspond to actual
ranking transitions.
\end{proposition}
\begin{proof}
Since the normalization factor is common to every coordinate, it cancels
when imposing
$H_\infty(P)=-\log\max_iP_i.$
Taking logarithms yields a linear equation in \(S\), whose unique solution
is the expression above. If the denominator vanishes, the ratio
\(S\) is constant and no isolated crossing occurs.
\end{proof}
\paragraph{Regularity Regions.}
Let
$T=\{t_{ij}\in(0,1)\},$
and order the distinct transition times
$0<\tau_1<\cdots<\tau_m<1.$
The complement
$[0,1]\setminus T$
is the union of finitely many open intervals.
On each interval the ordering of posterior probabilities is constant.
\section{Piecewise Smoothness}\label{ch:sec:piecewise-smoothness}
The previous statement alone is not sufficient: one must also assume
that the security functional is defined by smooth expressions once the
ordering is fixed.
\begin{theorem}[Piecewise Smoothness]
\label{ch:thm:17-1}
Let \(S\) be a security functional satisfying:
\begin{enumerate}
\item \(S\) depends on the posterior only through finitely many
    comparisons, maxima, minima, or rankings;
\item once the ordering of probabilities is fixed, \(S\) is given by a
    \(C^{k}\) expression (\(k\ge1\)).
\end{enumerate}
Then \(S\) is \(C^{k}\) on every regularity interval determined by the
transition times.
\end{theorem}
\begin{proof}
On a regularity interval the ordering of the posterior probabilities is constant. Consequently every comparison, maximum, minimum, or ranking operator in the definition of \(S\) reduces to a fixed selection of indices. Under assumption (2), the resulting expression is a \(C^{k}\) function of the posterior coordinates; since the WIS geodesic \(P(t)\) is smooth, the composition is \(C^{k}\) in \(t\) on the interval.
\end{proof}
\subsection{Local Geometric Bounds}\label{ch:cor:17-1}
On every regularity interval the hypotheses of Chapter~\ref{ch:sec:information-functionals-on-the-wis-manifold} apply.
\begin{corollary}[Local Stability]

If the WIS gradient of the local representative is bounded,
$\|S(P)-S(Q)\| \le L_S \sqrt{\mathcal D(P,Q)}$
for every pair of points contained in the same regularity interval.
If the local representative is \(C^{2}\) and has a stationary point,
\[ \|S(P)-S(P^{\star})\| \le C_S
\mathcal D(P,P^{\star}). \]
These estimates need not hold across ranking transitions without
additional analysis.
\end{corollary}

\section{Examples}\label{ch:sec:examples}

\paragraph{Min-Entropy.}
\(H_\infty(P)=-\log\max_iP_i.\)
When the maximizing index is unique, the functional reduces locally to
\(-\log P_{i^{\star}},\) which is smooth.

\paragraph{Bayes Vulnerability.}
\(V(P)=\max_iP_i.\)
Again, smoothness holds whenever the maximizing index is fixed.

\paragraph{Guessing Entropy.}
Guessing entropy depends on a permutation ordering the probabilities.
Once that permutation is fixed, the functional becomes a fixed weighted
sum of the probabilities and is therefore smooth.
\paragraph{Geometric Interpretation.}
The WIS manifold remains smooth throughout.
The only source of non-differentiability is the combinatorial change in
the ordering of posterior probabilities.
The transition times identify exactly where the analytical expression of
a rank-based functional changes.


\section{Discussion}\label{ch:sec:discussion-piecewise_smooth}
This chapter extends the WIS framework from globally smooth information
functionals to a large class of piecewise-smooth security measures.
The essential ingredients are:
\begin{enumerate}
\item WIS geodesics describe posterior evolution.
\item Ranking transitions are explicitly computable.
\item Between transitions the functional is smooth.
\item Universal geometric bounds apply locally on each regularity
    interval.
\end{enumerate}
The final chapter combines the smooth and piecewise-smooth theories into
a unified collection of security bounds.

